\documentclass[10pt, letterpaper]{article}

% Packages:
\usepackage[
    ignoreheadfoot, % set margins without considering header and footer
    top=2 cm, % seperation between body and page edge from the top
    bottom=2 cm, % seperation between body and page edge from the bottom
    left=2 cm, % seperation between body and page edge from the left
    right=2 cm, % seperation between body and page edge from the right
    footskip=1.0 cm, % seperation between body and footer
    % showframe % for debugging 
]{geometry} % for adjusting page geometry
\usepackage{titlesec} % for customizing section titles
\usepackage{tabularx} % for making tables with fixed width columns
\usepackage{array} % tabularx requires this
\usepackage[dvipsnames]{xcolor} % for coloring text
\definecolor{primaryColor}{RGB}{0, 0, 0} % define primary color
\usepackage{enumitem} % for customizing lists
\usepackage{amsmath, amssymb} % for math
\usepackage[
    pdftitle={Akshay Popat CV},
    pdfauthor={Akshay Popat},
    pdfcreator={LaTeX with RenderCV},
    colorlinks=true,
    urlcolor=primaryColor
]{hyperref} % for links, metadata and bookmarks
\usepackage[pscoord]{eso-pic} % for floating text on the page
\usepackage{calc} % for calculating lengths
\usepackage{bookmark} % for bookmarks
\usepackage{lastpage} % for getting the total number of pages
\usepackage{changepage} % for one column entries (adjustwidth environment)
\usepackage{paracol} % for two and three column entries
\usepackage{ifthen} % for conditional statements
\usepackage{needspace} % for avoiding page brake right after the section title
\usepackage{iftex} % check if engine is pdflatex, xetex or luatex

% Ensure that generate pdf is machine readable/ATS parsable:
\ifPDFTeX
    \input{glyphtounicode}
    \usepackage{lmodern}
\fi

% Some settings:
\raggedright
\AtBeginEnvironment{adjustwidth}{\partopsep0pt} % remove space before adjustwidth environment
\pagestyle{empty} % no header or footer
\setcounter{secnumdepth}{0} % no section numbering
\setlength{\parindent}{0pt} % no indentation
\setlength{\topskip}{0pt} % no top skip
\setlength{\columnsep}{0.15cm} % set column separation
\pagenumbering{gobble} % no page numbering



\titleformat{\section}{\needspace{4\baselineskip}\bfseries\large}{}{0pt}{}[\vspace{1pt}\titlerule]

\titlespacing{\section}{
    % left space:
    -1pt
}{
    % top space:
    0.3 cm
}{
    % bottom space:
    0.2 cm
} % section title spacing

\renewcommand\labelitemi{$\vcenter{\hbox{\small$\bullet$}}$} 
\newenvironment{highlights}{
    \begin{itemize}[
        topsep=0.10 cm,
        parsep=0.10 cm,
        partopsep=0pt,
        itemsep=0pt,
        leftmargin=0 cm + 10pt
    ]
}{
    \end{itemize}
} 


\newenvironment{highlightsforbulletentries}{
    \begin{itemize}[
        topsep=0.10 cm,
        parsep=0.10 cm,
        partopsep=0pt,
        itemsep=0pt,
        leftmargin=10pt
    ]
}{
    \end{itemize}
} 

\newenvironment{onecolentry}{
    \begin{adjustwidth}{
        0 cm + 0.00001 cm
    }{
        0 cm + 0.00001 cm
    }
}{
    \end{adjustwidth}
} 

\newenvironment{twocolentry}[2][]{
    \onecolentry
    \def\secondColumn{#2}
    \setcolumnwidth{\fill, 4.5 cm}
    \begin{paracol}{2}
}{
    \switchcolumn \raggedleft \secondColumn
    \end{paracol}
    \endonecolentry
} % new environment for two column entries

\newenvironment{threecolentry}[3][]{
    \onecolentry
    \def\thirdColumn{#3}
    \setcolumnwidth{, \fill, 4.5 cm}
    \begin{paracol}{3}
    {\raggedright #2} \switchcolumn
}{
    \switchcolumn \raggedleft \thirdColumn
    \end{paracol}
    \endonecolentry
} 

\newenvironment{header}{
    \setlength{\topsep}{0pt}\par\kern\topsep\centering\linespread{1.5}
}{
    \par\kern\topsep
}
\let\hrefWithoutArrow\href


\begin{document}
    \newcommand{\AND}{\unskip
        \cleaders\copy\ANDbox\hskip\wd\ANDbox
        \ignorespaces
    }
    \newsavebox\ANDbox
    \sbox\ANDbox{$|$}

    \begin{header}
        \fontsize{25 pt}{25 pt}\selectfont Akshay Popat

        \vspace{5 pt}

        \normalsize

        \mbox{\hrefWithoutArrow{mailto:akshay.popat23@imperial.ac.uk}{akshay.popat23@imperial.ac.uk}}%

    \end{header}

    \vspace{5 pt - 0.3 cm}


    %\section{Welcome to RenderCV!}




    %    \begin{onecolentry}
    %        \href{https://rendercv.com}{RenderCV} is a LaTeX-based CV/resume version-control and maintenance app. It allows you to create a high-quality CV or resume as a PDF file from a YAML file, with \textbf{Markdown syntax support} and \textbf{complete control over the LaTeX code}.
    %    \end{onecolentry}

    %    \vspace{0.2 cm}

    %    \begin{onecolentry}
    %        The boilerplate content was inspired by \href{https://github.com/dnl-blkv/mcdowell-cv}{Gayle McDowell}.
    %    \end{onecolentry}



    %\section{Quick Guide}

    %\begin{onecolentry}
    %    \begin{highlightsforbulletentries}


    %    \item Each section title is arbitrary and each section contains a list of entries.

    %    \item There are 7 unique entry types: \textit{BulletEntry}, \textit{TextEntry}, \textit{EducationEntry}, \textit{ExperienceEntry}, \textit{NormalEntry}, \textit{PublicationEntry}, and \textit{OneLineEntry}.

    %    \item Select a section title, pick an entry type, and start writing your section!

    %    \item \href{https://docs.rendercv.com/user_guide/}{Here}, you can find a comprehensive user guide for RenderCV.


    %    \end{highlightsforbulletentries}
    %\end{onecolentry}

    \section{INTERESTS}
    
        \begin{onecolentry}
            Homotopy theory, Algebraic and geometric topology
        \end{onecolentry}

        \vspace{0.2 cm}

    \section{EDUCATION}



        
        \begin{twocolentry}{
            10/2023 – present
        }
            \textbf{Imperial College London}, Mathematics
        \end{twocolentry}

        \vspace{0.10 cm}
        \begin{onecolentry}
            \begin{highlights}
                \item 3rd year average: 80.7
                \item Currently taking: Manifolds, Algebra 4, Commutative algebra
                \item Select 3rd year modules: Algebraic topology (91), Algebraic number theory (91), Galois theory (87), Group representation theory (86)
            \end{highlights}
        \end{onecolentry}

        \vspace{0.2 cm}

        \begin{twocolentry}{
            09/2015 – 06/2022
        }
            \textbf{Tiffin School}
        \end{twocolentry}

        \vspace{0.10 cm}
        \begin{onecolentry}
            \begin{highlights}
                \item A-Levels: A*A*A* (Mathematics, Further mathematics, Physics)
            \end{highlights}
        \end{onecolentry}

        \vspace{0.2 cm}




    %\section{Experience}




    %    \begin{twocolentry}{
    %        June 2005 – Aug 2007
    %    }
    %        \textbf{Software Engineer}, Apple -- Cupertino, CA\end{twocolentry}

    %    \vspace{0.10 cm}
    %    \begin{onecolentry}
    %        \begin{highlights}
    %            \item Reduced time to render user buddy lists by 75\% by implementing a prediction algorithm
    %            \item Integrated iChat with Spotlight Search by creating a tool to extract metadata from saved chat transcripts and provide metadata to a system-wide search database
    %            \item Redesigned chat file format and implemented backward compatibility for search
    %        \end{highlights}
    %    \end{onecolentry}


    %    \vspace{0.2 cm}

    %    \begin{twocolentry}{
    %        June 2003 – Aug 2003
    %    }
    %        \textbf{Software Engineer Intern}, Microsoft -- Redmond, WA\end{twocolentry}

    %    \vspace{0.10 cm}
    %    \begin{onecolentry}
    %        \begin{highlights}
    %            \item Designed a UI for the VS open file switcher (Ctrl-Tab) and extended it to tool windows
    %            \item Created a service to provide gradient across VS and VS add-ins, optimizing its performance via caching
    %            \item Built an app to compute the similarity of all methods in a codebase, reducing the time from $\mathcal{O}(n^2)$ to $\mathcal{O}(n \log n)$
    %            \item Created a test case generation tool that creates random XML docs from XML Schema
    %            \item Automated the extraction and processing of large datasets from legacy systems using SQL and Perl scripts
    %        \end{highlights}
    %    \end{onecolentry}



    
    %\section{Publications}




    %    \begin{samepage}
    %        \begin{twocolentry}{
    %            Jan 2004
    %        }
    %            \textbf{3D Finite Element Analysis of No-Insulation Coils}
    %        \end{twocolentry}

    %        \vspace{0.10 cm}

    %        \begin{onecolentry}
    %            \mbox{Frodo Baggins}, \mbox{\textbf{\textit{John Doe}}}, \mbox{Samwise Gamgee}

    %            \vspace{0.10 cm}

    %    \href{https://doi.org/10.1109/TASC.2023.3340648}{10.1109/TASC.2023.3340648}
    %    \end{onecolentry}
    %    \end{samepage}


    
    \section{PROJECTS}

        \begin{twocolentry}{
            10/2026 -- present
        }
            \textbf{Hitchin systems and their quantization}
        \end{twocolentry}

        \vspace{0.10 cm}
        \begin{onecolentry}
            \begin{highlights}
                \item 4th year project supervised by Professor Yankı Lekili
            \end{highlights}
        \end{onecolentry}

        \vspace{0.2 cm}


        \begin{twocolentry}{
            09/2026
        }
            \textbf{EPSRC Vacation Internship -- Principal bundles}
        \end{twocolentry}

        \vspace{0.10 cm}
        \begin{onecolentry}
            \begin{highlights}
                \item Selected for an EPSRC-funded internship supervised by Dr Marie-Amelie Lawn
                \item Proved the equivalence between rank \(n\) real vector bundles and principal \(GL_n(\mathbb{R})\)-bundles, and studied their classification via the infinite Grassmannian \(Gr_n(\mathbb{R}^\infty)\)
                \item Explicitly constructed examples of spin structures
                \item Proved that an \(SO(n)\)-bundle over a manifold \(M\) has a spin structure if and only if the second Stiefel-Whitney class vanishes, in which case the spin structures are classified by \(H^1(M; \mathbb{Z}_2)\)
            \end{highlights}
        \end{onecolentry}

        \vspace{0.2 cm}

        \begin{twocolentry}{
            06/2025
        }
            \textbf{Multilinear Kakeya Estimates}
        \end{twocolentry}

        \vspace{0.10 cm}
        \begin{onecolentry}
            \begin{highlights}
                \item 2nd year group project supervised by Dr Zhiyuan Zhang
                \item Proved the Borsuk-Ulam theorem and learned of its application in proving a weak Bennett-Carbery-Tao multilinear Kakeya estimate
            \end{highlights}
        \end{onecolentry}


        \vspace{0.2 cm}

    \section{TECHNICAL SKILLS}

        \begin{onecolentry}{
            
        }
            \LaTeX, Python
        \end{onecolentry}

        \vspace{0.2 cm}

    \section{TUTORING}

        \begin{twocolentry}{
            01/2026 -- present
        }
            \textbf{We Solve Problems -- KCL Maths Circle Tutor}
        \end{twocolentry}

        \vspace{0.10 cm}
        \begin{onecolentry}
            \begin{highlights}
                \item Help students through material and problems
            \end{highlights}
        \end{onecolentry}

        \vspace{0.2 cm}

    \section{ADDITIONAL}

        \begin{twocolentry}{
            10/2024 -- present
        }
            \textbf{Imperial College MathsWiki}
        \end{twocolentry}

        \vspace{0.10 cm}
        \begin{onecolentry}
            \begin{highlights}
                \item Contributing to making notes for the undergraduate courses at Imperial
            \end{highlights}
        \end{onecolentry}

        \vspace{0.2 cm}

                \begin{twocolentry}{
            07/2022 -- present
        }
            \textbf{Japanese}
        \end{twocolentry}

        \vspace{0.10 cm}
        \begin{onecolentry}
            \begin{highlights}
                \item Self-studying Japanese
                \item N1 level
            \end{highlights}
        \end{onecolentry}

\end{document}
